\documentclass[11pt]{article}
\usepackage[T1]{fontenc}
\usepackage{textcomp}
\usepackage[margin=0.75in]{geometry}
\usepackage{amsmath,amssymb,amsthm}
\usepackage{array,booktabs}
\usepackage{hyperref}
\setlength{\parindent}{0pt}
\newtheorem{theorem}{Theorem}
\theoremstyle{definition}
\newtheorem{definition}{Definition}
\title{Flow Proof Artifact: prop-29-equal-arcs-equal-chords}
\author{Generated by \texttt{flow doc proof}}
\date{Flow Proof Book}
\begin{document}
\maketitle
In equal circles equal circumferences are subtended by equal straight lines.\\[0.5em]
\textbf{Source.} euclid — Elements, Book III, Proposition 29\\[0.5em]
\bigskip
\noindent{\large \textbf{Derived fact 1.} \textit{In equal circles equal circumferences are subtended by equal straight lines} \hfill the Euclidean plane \cdot Euclid Book III \cdot Proposition 29: in equal circles equal arcs subtend equal chords}\par
\smallskip\noindent\textit{Coordinate: the Euclidean plane \cdot Euclid Book III \cdot Proposition 29: in equal circles equal arcs subtend equal chords}\par
\smallskip\noindent\textit{Source: euclid — Elements, Book III, Proposition 29}\par
\smallskip\noindent\textit{Built on: proposition 28: in equal circles equal chords cut off equal arcs, for Euclid Book III on the Euclidean plane}\par
\medskip
\noindent\textbf{Goal.} In equal circles equal circumferences are subtended by equal straight lines\par
\noindent$\displaystyle \text{equal arcs subtend equal chords}$\par
\medskip
\noindent
\begin{tabular}{@{} >{\raggedright\arraybackslash}p{0.06\textwidth} >{\raggedright\arraybackslash}p{0.47\textwidth} >{\raggedleft\arraybackslash}p{0.41\textwidth} @{}}
\textbf{\#} & \textbf{Proof} & \textbf{Mathematics} \\
\midrule
\textbf{1.} & We prove that proposition 29: in equal circles equal arcs subtend equal chords for Euclid Book III the Euclidean plane. & \\[0.45em]
\textbf{2.} & Let arc1 = equal\_arc. & \\[0.45em]
\textbf{3.} & Let arc2 = equal\_arc. & \\[0.45em]
\textbf{4.} & Let chord1 = chord\_subtending\_arc1. & \\[0.45em]
\textbf{5.} & From step 2, step 3, and step 4, we can deduce that chord1 equals chord2. & $\displaystyle chord1 = chord2$ \\[0.45em]
\textbf{6.} & We invoke the derived fact governing Euclid Book III the Euclidean plane: proposition 28: in equal circles equal chords cut off equal arcs, for Euclid Book III on the Euclidean plane. & \\[0.45em]
\textbf{7.} & From step 6, this implies equal arcs subtend equal chords. Hence proven. & $\displaystyle \text{equal arcs subtend equal chords}$ \\[0.45em]
\bottomrule
\end{tabular}
\label{thm:Geometry-Euclid Book III-proposition_29_in_equal_circles_equal_arcs_subtend_equal_chords}
\par\medskip\hrule
\end{document}
